\documentclass[10pt,a4paper]{article}
\usepackage[a4paper,margin=18mm]{geometry}
\usepackage[T1]{fontenc}
\usepackage[utf8]{inputenc}
\usepackage{lmodern,microtype,amsmath,amssymb,booktabs,tabularx,array,graphicx,xcolor,enumitem,fancyhdr}
\usepackage[font=small,labelfont=bf]{caption}
\usepackage{titlesec}
\usepackage{xurl}
\urlstyle{tt}
\usepackage[colorlinks=true,linkcolor=teal,urlcolor=teal]{hyperref}
\definecolor{teal}{HTML}{086C7B}
\definecolor{amber}{HTML}{A64B00}
\newcommand{\sci}[2]{\ensuremath{#1\times10^{#2}}}
\newcommand{\file}[1]{{\footnotesize\nolinkurl{#1}}}
\newcommand{\status}[1]{\textbf{\textcolor{amber}{#1}}}
\newcommand{\team}{Maybe Once}
\graphicspath{{../pitch_outputs/}{figures/}{./}}
\newcommand{\plot}[2]{\IfFileExists{../pitch_outputs/#1}{\includegraphics[width=#2]{../pitch_outputs/#1}}{\IfFileExists{figures/#1}{\includegraphics[width=#2]{figures/#1}}{\IfFileExists{#1}{\includegraphics[width=#2]{#1}}{\fbox{\parbox[c][30mm][c]{\dimexpr#2-8mm\relax}{\centering Upload \file{#1} to Overleaf}}}}}}
\setlength{\parindent}{0pt}\setlength{\parskip}{4pt}
\setlength{\tabcolsep}{4pt}\renewcommand{\arraystretch}{1.13}
\newcolumntype{L}[1]{>{\raggedright\arraybackslash}p{#1}}
\renewcommand{\tabularxcolumn}[1]{>{\raggedright\arraybackslash}p{#1}}
\setlist{nosep,leftmargin=1.2em}
\titlespacing*{\section}{0pt}{7pt}{4pt}
\titleformat{\section}{\large\bfseries\color{teal}}{\thesection}{0.5em}{}
\pagestyle{fancy}\fancyhf{}\fancyhead[L]{\small OctoEarth Simulator}
\fancyhead[R]{\small Technical report / 25 October 2026}\fancyfoot[C]{\thepage}
\setlength{\headheight}{14pt}
\emergencystretch=2em
\begin{document}
% Eight deliberately separated content pages. Verify the compiled page count in Overleaf.
\begin{center}
{\LARGE\bfseries Moving the Earth--Moon System\\with a Remote Gravity Tractor}\\[5pt]
{\large Solar-powered propulsion, conservation and feasibility limits}\\[5pt]
\team\quad |\quad Planetary Propulsion Simulation Hackathon\\
\small NIT Trichy
\end{center}
\section{Executive result and challenge alignment}
\textbf{We demonstrate conditional orbital migration and identify a failed structural requirement; we do not demonstrate a successful planetary relocation.} A massive tractor pulls Earth and Moon through reciprocal gravity. Solar-powered, canted reaction-mass jets maintain the tractor's commanded location; exhaust carries the reaction momentum. A collector supplies energy and a mirror attempts to compensate for declining sunlight.

The selected design's preserved one-hour-step recording ends after 5460.78 years at an Earth--Moon barycentre (EMB) semi-major axis of 1.76381 AU. Its minimum structural safety factor is 0.13897, below the required 1.5. A different, earlier layout crossed 2.7 AU after 8742.18 years but failed safety and engineering checks. These are distinct designs and evidence sets.

\par\smallskip
{\small\noindent\begin{tabularx}{\linewidth}{@{}L{33mm}X L{24mm}@{}}\toprule
Guide win condition & Best available evidence & Assessment\\\midrule
Target within 10,000 yr & Selected arrival not observed; original crossed at 8742.18 yr & Incomplete selected\\
Planet remains safe & Proxy checks only; original has shading, tide, flux and plume flags & Not established\\
Every structural SF $\ge1.5$ & Selected minimum 0.13895; original minimum 0.09210 & \status{FAIL}\\
Applicable benchmarks pass & B1--B10 recorded passing; newer fast suite 35 passed, 0 failed, 1 skipped & Tested scope only\\
Conservation and reactions & Source/sink ledger residuals below $10^{-9}$; original terminal ledger available & Numerical closure\\
No unresolved non-enabler flags & Four remain in selected prefixes & \status{FAIL}\\\bottomrule
\end{tabularx}\par}
\smallskip


\textbf{Contribution.} A Unity-independent double-precision C\# core couples 13-body gravity, finite booked station-keeping impulses, solar supply, material consumption, radiation, climate and engineering monitors. Its main scientific result is the separation of attainable motion in an assumed model from unsupported engineering feasibility.

\textbf{Evidence scope.} The selected integrations are interrupted recordings, not completed transfers. No complete restart history or terminal selected report exists. Thresholds in the implementation are explicitly placeholders; passing them does not certify planetary safety.


\newpage
\section{Concept, architecture and declared assumptions}
The tractor is maintained ahead of Earth approximately along the EMB prograde direction, with a selected $5^\circ$ Sunward bias. Adding prograde acceleration increases orbital energy and semi-major axis; final circular orbital speed decreases as the orbit expands. The tractor's backward gravitational attraction is countered by jets. Ejected mass and incoming replacement material are explicitly booked, so this is an open system with named momentum partners.

\par\smallskip
{\small\noindent\begin{tabularx}{\linewidth}{@{}Xrr@{}}\toprule
Parameter & Selected design & Original layout\\\midrule
Initial tractor mass (kg) & $1.0098073\times10^{19}$ & $1.6104890\times10^{18}$\\
Earth-centre separation (km) & 125,000 & 50,000\\
Exhaust speed / cant angle & 300 km/s / $25^\circ$ & Same\\
Sunward bias / sizing margin & $5^\circ$ / 1.15 & $0^\circ$ / 1.15\\
Collector placement & Co-moving, phase $180^\circ$ & Fixed 1 AU, phase $0^\circ$\\
Collector inclination & $10^\circ$ & $10^\circ$\\
Post-target policy & Depart (unexercised) & Powered parking\\\bottomrule
\end{tabularx}\par}
\smallskip


At assumed density 3000 kg/m$^3$, the selected spherical tractor radius is 92.97 km. Its acquisition and placement are assumed, not simulated.

\textbf{Exactly three declared enablers.}
\par\smallskip
{\small\noindent\begin{tabularx}{\linewidth}{@{}lXX@{}}\toprule
ID & Bounded declaration & Consequence\\\midrule
E1 & Material at consumer velocity; $3\times10^{23}$ kg cap, $10^{13}$ kg/s delivery & Replaces mining/transport; books mass, momentum and energy through EnablerSource\\
E2 & Ideal transmission, conversion, thrusters and station keeping; jet efficiency 1, $v_e=300$ km/s & Hardware idealization; finite solar supply and material still required\\
E3 & Initially placed selected tractor, $1.0098073\times10^{19}$ kg & Replaces tractor acquisition/placement\\\bottomrule
\end{tabularx}\par}
\smallskip

No hypothetical energy source is declared. Giant focusing optics remain unresolved outside E1--E3; organizer acceptance has not been inferred.

\textbf{Software and data flow.} Configuration and approximate J2000 data $\rightarrow$ initial conditions and tractor sizing $\rightarrow$ gravity, control impulses, resources and radiation $\rightarrow$ climate and monitors $\rightarrow$ CSV, JSON and ledger $\rightarrow$ Unity playback. The 13 bodies are Sun, eight planets, Moon, tractor, collector and mirror. Ejecta are ledger parcels, not propagated N-body particles.

\textbf{Model labels.} S (simplified): Newtonian ephemeris, climate, radiation, local structure. E (enabler-dependent): delivery, ideal hardware and initial placement. P (placeholder/unvalidated): safety thresholds, strengths, focusing optics and global thermal structure. Numerical convergence validates specified algorithms, not these engineering assumptions.

\textbf{Implementation.} Numerical sources reside in \file{Assets/Scripts/Core/}; the headless entry point is \file{Tools/OctoEarth.Cli/Program.cs}. The shared project links those same core sources. Unity 6000.3.0f1 renders recorded mission state; a separate baseline mode can integrate live. Display coordinates never feed back into mission physics.

\textbf{Why this alternative?} Earth's own field and internal structural forces cannot provide external propulsion; a uniform-field closed current loop has zero net force. The guide's 60 N wire estimate is far below the required order $10^{17}$ N. This motivates an explicitly external reaction mechanism; it does not prove that a gravity tractor is uniquely feasible.
\newpage
\section{Governing calculations and physical models}
For Newtonian bodies and a heliocentric orbital diagnostic,
\begin{align}
\ddot{\mathbf r}_i&=\sum_{j\ne i}Gm_j\frac{\mathbf r_j-\mathbf r_i}{|\mathbf r_j-\mathbf r_i|^3},&
\varepsilon&=\frac{v^2}{2}-\frac{\mu}{r}=-\frac{\mu}{2a},\\
\frac{da}{dt}&=\frac{2a^2}{\mu}\mathbf v\cdot\mathbf a_{\rm push},&
\Delta E&=\frac{\mu M}{2}\left(\frac1{a_1}-\frac1{a_2}\right).
\end{align}
The circular-endpoint sizing budget uses $M=6.045\times10^{24}$ kg (a rounded Earth--Moon mass), $a_1=1$ AU, $a_2=2.7$ AU and $\tau=10^4$ Julian years. Dynamical body masses use their separate constants, not this rounded pair mass.
\[
\Delta v=\sqrt{\mu/a_1}-\sqrt{\mu/a_2},\quad p=M\Delta v,\quad
\bar F=p/\tau,\quad\bar P=\Delta E/\tau.
\]
\par\smallskip
{\small\noindent\begin{tabularx}{\linewidth}{@{}XrXr@{}}\toprule
Quantity & Computed sizing & Quantity & Computed sizing\\\midrule
$\Delta E$ & $1.6883\times10^{33}$ J & $\Delta v$ & 11.6583 km/s\\
$p$ & $7.0474\times10^{28}$ N s & $\bar F$ & $2.2332\times10^{17}$ N\\
$\bar P$ & $5.3498\times10^{21}$ W & $\Delta L$ & $1.7324\times10^{40}$ kg m$^2$/s\\\bottomrule
\end{tabularx}\par}
\smallskip

These are analytic planning outputs, not a measured transfer. $\bar P$ is orbital energy gain per time, not total jet electrical power.

For symmetric canted jets, lateral impulses cancel while net thrust $F$ gives
\[
\dot m=\frac{F}{v_e\cos\varphi},\qquad P_{\rm jet}=\frac{Fv_e}{2\cos\varphi},\qquad
A_{\rm col}=\frac{(1+f_{\rm reserve})P_{\rm req}}{S\eta_c-p_{\rm sk}}.
\]
The selected initial collector area is $1.09516\times10^{20}$ m$^2$; areal density 0.01 kg/m$^2$ implies $1.09516\times10^{18}$ kg collector mass. Efficiency is 0.30, reflectivity 0.10, reserve 0.05; area is capped at $10^{21}$ m$^2$. Mirror reflectivity is 0.95, nominal distance $10^9$ m, optical error 0.001 rad and area cap $10^{17}$ m$^2$.

\textbf{Climate.} The implemented global energy balance is
\[
C\dot T=\frac{(1-\alpha(T))S_{\rm vis}+S_{\rm IR}}4-\mathrm{OLR}(T),\quad
C=1000(4186)H,\quad \alpha=0.30+0.32f_{\rm ice}.
\]
$H=100$ m; ice fraction changes linearly from 0 at 273.15 K to 1 at 263.15 K. Above 233.15 K, $\mathrm{OLR}=203.3+2.09(T-273.15)$ W/m$^2$; below it a continuous emissive $T^4$ branch is used. Infrared is absorbed separately from visible albedo. This is an illustrative model without clouds, carbon cycle or spatial climate. The unheated radiative equilibrium estimate is 254.04 K at 1 AU and 154.60 K at 2.7 AU; neither is a full surface-climate prediction.

\textbf{Structure.} Local beams use $\sigma=N/A\pm Mc/I$, $P_{\rm cr}=\pi^2EI/\ell^2$ and $\mathrm{SF}=\sigma_{\rm allow}/\sigma$. The governing tractor proxy is $\sigma_{\rm rock}=\rho(2Ra+GM_T/R)$. It has no resolved stress tensor, confinement or interface mechanics; its 5 MPa allowable is a generic unvalidated input, not a measured mode-specific strength.
\newpage
\section{Numerical method and validation}
State uses C\# \texttt{double}, custom \texttt{Vector3d} and compensated Kahan summation. Initial conditions are approximate J2000 elements with an ideal circular, coplanar Moon, not a validated long-term ephemeris. Yoshida4 advances conservative gravity; positive-time external half-kicks make the forced mission approximately second order. Prescribed infrastructure paths are enforced by finite booked impulses with a 0.01 m endpoint solver tolerance and at most 16 iterations; this is not autonomous guidance validation.

Selected timesteps are 3600 and 1800 s. Climate uses RK4 substeps no longer than one day; finite-disc IR is refreshed after six hours or area change. Recording is nominally monthly plus checkpoints. Encke deviation tracking is tested separately and is not the mission integrator.

\par\smallskip
{\small\noindent\begin{tabularx}{\linewidth}{@{}r L{49mm} X@{}}\toprule
B & Benchmark & Recorded measurement / implemented tolerance\\\midrule
1 & 10 A, 2 m $x$ wire, 0.5 T $y$ field & $+z$ force 10 N; vector error 0, tolerance $10^{-8}$ N\\
2 & Closed loop, uniform field & Net force 0; tolerance $5\times10^{-8}$ N\\
3 & 1 kA long wire, radius 1 m & $2.000\times10^{-4}$ T; tolerance $2\times10^{-10}$ T\\
4 & Kepler $a=1$ AU, $e=0.3$, ten periods & Return error $2.51\times10^{-11}a$, energy $3.29\times10^{-13}$; limits $10^{-6}a$, $10^{-9}$\\
5 & Full solar system, no push, five years & Energy drift $9.57\times10^{-14}$; limit $10^{-9}$\\
6 & EMB sidereal period & 365.260245 d; $365.256\pm0.1$ d\\
7 & Hohmann 1 to 1.524 AU & 2.94606 + 2.64998 km/s, 258.915 d; limits 0.02 km/s, 2 d\\
8 & $10^{-8}$ m/s$^2$ prograde, one year & $\Delta a=3170.101$ km; 3170 km $\pm1\%$\\
9 & Self-weight cantilever & Moment 5000 N m, deflection 6.25 mm; 1\%\\
10 & Equilibrium, 1 AU, albedo 0.306 & 254.0386 K; $254\pm2$ K\\\bottomrule
\end{tabularx}\par}
\smallskip

All ten groups have recorded passing comparisons. B1--B3 test electromagnetic utilities, not mission propulsion. The table states implemented absolute tolerances; the guide's B1--B3 tolerance notation is not identical. Zero measured B1/B2 error does not justify silently relabelling code tolerances as the guide's values.

\textbf{Additional evidence.} Observed conservative orders: Yoshida 3.998, Verlet 2.003. The selected ten-year five-step study (7200/3600/1800/900/450 s) gives successive semi-major-axis orders 2.012, 2.003 and 2.001. A tiny-push test gives 0 m in floats and 37.32354 m in doubles versus 37.32357 m expected. These are bounded tests, not long-duration accuracy guarantees.

\textbf{Test provenance.} The newer fast EditMode suite reports 35 passed, 0 failed, one explicit 10,000-year test skipped (13.875 s). A subsequent focused viewer regression passed once. The old bridge summary lacking passed/skipped counts is superseded for this evidence. SDK/xUnit execution was unavailable; Unity-bundled Roslyn headless benchmarks ran successfully. No new long baseline was run for the newer audit.

\textbf{Gate interpretation.} G1 hover and G2 delivered-velocity/Moon-binding checks passed. G2 orbital-difference envelopes and X6 endpoint envelopes are diagnostic, not precision acceptance. G3 is a separate Sun/Earth/tractor toy: target times 8758.15 and 8710.79 yr differ by 47.36 yr (0.544\%). Its analytic-error ratio 4.13 is consistent with second-order behaviour but does not establish selected mission arrival convergence.
\newpage
\section{Recorded migration, resources and timestep sensitivity}
\par\smallskip
{\small\noindent\begin{tabularx}{\linewidth}{@{}Xrrr@{}}\toprule
Evidence & Last time (yr) & $a_{\rm EMB}$ (AU) & $e_{\rm EMB}$\\\midrule
Selected $h=3600$ s, interrupted & 5460.7804 & 1.7638141 & 0.0103894\\
Selected $h=1800$ s, interrupted & 2361.1500 & 1.2515532 & 0.0129082\\
Original layout, completed & 10000 & 2.6991380 & 0.0207656\\\bottomrule
\end{tabularx}\par}
\smallskip

The original crossed 2.7 AU at 8742.1759 yr then underwent powered parking. Its final value differs from the crossing value; it is not evidence of selected arrival. Its climate history predates corrected infrared coupling. Selected post-target departure has not been exercised. Semi-major axis is an orbital diagnostic, not instantaneous Earth--Sun distance.

\begin{center}\plot{selected_orbit.png}{0.47\textwidth}\hfill\plot{selected_propellant.png}{0.47\textwidth}
\captionof{figure}{Selected recorded orbital growth and cumulative ejection. Curves derive from preserved CSVs; the target is not reached within the available recording.}
\end{center}

\par\smallskip
{\small\noindent\begin{tabularx}{\linewidth}{@{}Xrr@{}}\toprule
Resource quantity at each prefix endpoint & $h=3600$ s & $h=1800$ s\\\midrule
E1 mass used (kg) & $1.63526\times10^{23}$ & $7.07299\times10^{22}$\\
Fraction of $3\times10^{23}$ kg cap & 54.5087\% & 23.5766\%\\
Derived jet kinetic energy (J) & $7.35852\times10^{33}$ & $3.18277\times10^{33}$\\\bottomrule
\end{tabularx}\par}
\smallskip

Different endpoint times explain different totals. E1 includes infrastructure and supplied propellant; E3 covers the initially placed tractor. Jet energy is $m_{\rm ejecta}v_e^2/2$. Available solar electricity is checked for every ejection under E2. Cumulative generation and unused-energy totals were not saved for interrupted selected runs; their utilization fraction cannot be recovered from the sparse thrust column.

\textbf{Partial-duration comparison only.} Over the shared 2361.15 yr, coarse scalar outputs are linearly aligned onto fine timestamps; exact shared checkpoints are retained separately. Maximum differences are $1.50724\times10^{-4}$ AU in $a$, $1.14160\times10^{-4}$ in eccentricity, $2.68626\times10^6$ m in Moon distance and $2.81310\times10^{-4}$ K in temperature. Monthly records cannot establish all-step extrema or resolve lunar phase accurately. Alignment corrections are diagnostics, not rigorous error bounds.

The final coarse-minus-fine ejecta difference is $-2.34548\times10^{19}$ kg, about $-0.03316\%$ of the fine total. It accumulates at nearly constant rate and is much larger than the $1.70828\times10^{15}$ kg maximum alignment correction. Ejection sums actual half-kick bookings from all three controlled bodies; the thrust scalar represents only the tractor's final kick. Thus a sparse thrust integral cannot explain away this difference. Its precise numerical cause remains unresolved.
\newpage
\section{Planetary environment, optics and structural failure}
All implementation safety thresholds are marked placeholders. Absence of a flag is a statement about these model checks, not proof of habitability. Neither prefix added a shading, plume, tide or flux flag beyond its initial four flags.

\par\smallskip
{\small\noindent\begin{tabularx}{\linewidth}{@{}L{34mm}L{28mm}X@{}}\toprule
Metric & Retained criterion & Selected evidence and limitation\\\midrule
Temperature & 278--298 K & Sampled 288--290.8872 K; simplified EBM\\
Tide / lunar tide & $\le0.01$ & Sampled maximum approximately 0.0044965; gradient proxy\\
Incoming flux & $\le1429.05$ W/m$^2$ & Sampled maximum approximately 1407.93 W/m$^2$\\
Solar-disc shading & $\le0.01$ & No prefix flag; original reaches full shading\\
Moon distance & 300,000--500,000 km & Sampled approximately 374,491--399,762 km\\
Moon binding & Negative specific energy; inside Hill radius & Cumulative bound true in both prefixes; no future guarantee\\
Plume exclusion / heating & $3R_E$ / 1 W/m$^2$ & No prefix flag; long-prefix terminal clearance extrema missing\\
Eccentricity & $\le0.1$ & Recorded endpoints 0.01039 / 0.01291\\
Atmosphere, spin, magnetic shielding, biological dose & Not implemented & No safety conclusion available\\\bottomrule
\end{tabularx}\par}
\smallskip


\textbf{Mirror dependence.} Thermal-only replays remove reflected mirror light while retaining recorded direct illumination, array infrared and trajectory. Daily RK4 substeps use monthly held forcing; infrastructure, expenditure and dynamics are not recomputed. Temperature first leaves 278--298 K near year 605.98 in both replays. Final temperatures are 179.13 K (longer coarse prefix) and 212.59 K (shorter fine prefix), with ice fraction 1. These different endpoints are not a same-time convergence comparison. The replay demonstrates conditional cooling, not a complete mission with optics rejected.

At 1 AU, the implemented spot geometry gives solar angular diameter 0.009301 rad and full interception only if
\[
\delta^2\le(2R_E/d_m)^2-\theta_\odot^2.
\]
For $d_m=10^9$ m, $\delta_{\max}=0.008709$ rad. At $2\times10^9$ and $5\times10^9$ m no real full-interception bound exists in this geometry even at zero error. This is an optical sensitivity, not proof of a buildable mirror or late-time pointing stability.

\textbf{Structural verdict: FAIL at the unchanged 5 MPa allowable.}
\par\smallskip
{\small\noindent\begin{tabularx}{\linewidth}{@{}Xrrr@{}}\toprule
Component/evidence & Stress (MPa) & Allowable (MPa) & SF\\\midrule
Selected coarse rock proxy & 35.9801 & 5 & 0.138966\\
Selected fine rock proxy & 35.9835 & 5 & 0.138953\\
Array local beam, one-year diagnostic & 31.9283 & 200 & 6.264\\
Mirror local beam, one-year diagnostic & 3.61274 & 200 & 55.36\\\bottomrule
\end{tabularx}\par}
\smallskip

The worst recorded proxy would require 53.9752 MPa for SF 1.5, a 10.795-fold increase; no such strength is adopted. Local beam results do not validate global structure or mounts. Small-specimen compression strengths cannot replace this generic scalar allowable: tension, shear, bending, joints, defects and scale require resolved analysis.

Four selected non-enabler flags remain: StructuralSafety, FocusingOpticsUnresolved, MaterialStrengthUnvalidated and ThermalStructuralModelIncomplete. The original has 14, additionally including tides, flux, shading, initially inward ejecta, plume heating and five plume-intersection flags.
\newpage
\section{Conservation and momentum record}
Gravity exchanges equal and opposite forces between simulated bodies. Jets transfer momentum to EjectaSink; replacement and construction material enter through EnablerSource. Solar pressure, mirror reflection and collector IR exchange momentum with PhotonSink. SolarArrayJetEnergySource books jet energy input. Negative source entries are accounting withdrawals from an external reservoir, not negative physical mass.

With tracked bodies $\mathcal B$ and external reservoirs $\mathcal R$, the bookkeeping checks have the form
\[
\Delta\mathbf P_{\mathcal B}+\Delta\mathbf P_{\mathcal R}\simeq0,\quad
\Delta\mathbf L_{\mathcal B}+\Delta\mathbf L_{\mathcal R}\simeq0,\quad
\Delta E_{\mathcal B}+\Delta E_{\mathcal R}\simeq0.
\]
Energy terms include supplied parcels and explicit jet input. Stored relative errors use initial gross body scales: $\sum m|v|$ for momentum, $\sum|\mathbf r\times m\mathbf v|$ for angular momentum and initial absolute mechanical energy for energy. They are not uniformly normalized by propulsion impulse. Separate force/reaction checks address impulse closure.

\par\smallskip
{\small\noindent\begin{tabularx}{\linewidth}{@{}Xrrr@{}}\toprule
Recorded cumulative maximum & Relative $P$ & Relative $L$ & Relative $E$\\\midrule
No-push baseline, 10,000 yr & $7.31\times10^{-13}$ & $1.07\times10^{-12}$ & $2.20\times10^{-12}$\\
Original propulsion, 10,000 yr & $6.23\times10^{-13}$ & $1.13\times10^{-12}$ & $7.56\times10^{-12}$\\
Selected coarse prefix & $6.77\times10^{-13}$ & $9.94\times10^{-13}$ & $6.73\times10^{-13}$\\
Selected fine prefix & $4.05\times10^{-13}$ & $4.89\times10^{-13}$ & $1.17\times10^{-12}$\\\bottomrule
\end{tabularx}\par}
\smallskip

All lie below the retained $10^{-9}$ threshold. The long baseline contains Sun, Earth and Moon; B5 separately tests the all-planet no-push system. Small residuals establish numerical/accounting closure within the model, not feasible hardware or a certified transfer.

\begin{center}\plot{selected_ledger.png}{0.62\textwidth}
\captionof{figure}{Selected recorded conservation history. Residuals are cumulative solver diagnostics with gross-scale normalization.}
\end{center}

\textbf{Complete original ledger.} The complete recorded ledger is \file{Results/mission_h3600_y10000/ledger.txt}, accompanied by its trajectory, final state and mission report. Applied and reaction accumulated vector-impulse magnitudes both equal $1.90583\times10^{26}$ N s. This is neither Earth-only delivered impulse nor a scalar sum of thrust magnitudes. Force closure is $3.63\times10^{-16}$, reaction closure $6.63\times10^{-16}$ and solar-sector closure $1.53\times10^{-16}$. The absolute final residual vector is not exported.

\textbf{Solar sector, original run only.} Incident collected sunlight is $4.70428\times10^{34}$ J; jet requirement is $1.33441\times10^{34}$ J, about 7.90 times the analytic orbital budget. Material consumption is $2.96537\times10^{23}$ kg. Reported sourceFraction = 1 is a clamped electrical-used/jet-required ratio, not source-to-orbit efficiency. Selected interrupted runs lack terminal generation totals and full terminal ledgers. Their residual records must not be presented as complete selected mission conservation records.
\newpage
\section{Assumptions, reproducibility and evidence}
\par\smallskip
{\small\noindent\begin{tabularx}{\linewidth}{@{}L{30mm}X@{}}\toprule
Unresolved assumption & Consequence if unsupported / next evidence needed\\\midrule
Giant focusing optics & Solar compensation conditional; obtain organizer interpretation and specify aperture, pointing and global stability\\
Material strengths and mounts & Structural proxy fails; resolve stress modes, interfaces and realistic scale before claiming a feasible tractor\\
Local thermal/structural proxies & Global array/mirror deformation and stability unvalidated; require coupled analysis\\
E2 ideal hardware and paths & Transmission, losses and autonomous control not demonstrated; recompute coupled energy/heat/guidance if rejected\\
Ballistic ejecta clearance & Launch cones do not prove future escape; propagate parcels and test return/collision risk\\
Global EBM and placeholders & Cannot certify habitability; justify thresholds and add atmospheric, spin, magnetic and radiation models\\
Selected arrival and departure & Unobserved; matched full-duration $h/h2$ runs required after feasibility review\\\bottomrule
\end{tabularx}\par}
\smallskip


\textbf{Reproducibility.} From the project root, the documented fallback build is \file{Tools/Build-HeadlessWithoutSdk.ps1}; it uses Unity's bundled compiler and installed .NET runtime. SDK/MSBuild and xUnit execution were not verified. The existing audited executable accepts:
{\footnotesize\begin{verbatim}
dotnet pitch_outputs/build/OctoEarth.Cli.dll bench
dotnet pitch_outputs/build/OctoEarth.Cli.dll mission --config
  Config/mission_design_candidate.json --dt 3600 --years 10000
  --output NEW_OUTPUT_DIRECTORY
\end{verbatim}}
The mission command is shown across lines for readability and is entered as one line. It starts from initial conditions. A matched fine integration uses 1800 s with identical physical configuration and a distinct output directory. Selected full-duration results remain unavailable.

\par\medskip
\noindent\textbf{Evidence index}\par
{\footnotesize\raggedright
Paths below are relative to the project root.\par
\textbf{Participant guide:} Planetary Propulsion Simulation Hackathon Participant Guide, sections 1, 6, 7 and 9.\par
\textbf{Selected design:} \file{Docs/Selected_Design_Prefix_Report.md}; \file{Results/partial_convergence.json}. Raw recordings: \file{Results/design_audit/full_h3600/run_mission.csv} and \file{Results/design_audit/full_h1800/run_mission.csv}.\par
\textbf{Presentation and replay evidence:} \file{Docs/Playback_Presentation_Audit.md}; \file{Results/moon_report.json}; \file{Results/no_mirror_replay_h3600.csv}; \file{Results/no_mirror_replay_h1800.csv}; \file{Results/mirror_optics_sensitivity.csv}; \file{Results/presentation_editmode_tests.json}.\par
\textbf{Original mission:} \file{Results/mission_h3600_y10000/} contains \file{mission_report.json}, \file{ledger.txt}, \file{final_state.json}, \file{run_mission.csv} and \file{enablers.json}.\par
\textbf{Baseline and validation:} \file{Results/baseline_validation.json}; \file{Results/run_baseline_10000.csv}; \file{pitch_outputs/Results/benchmarks.json}; \file{Results/design_audit/order/}.\par
\textbf{Implementation:} \file{Assets/Scripts/Core/}; \file{Config/mission_design_candidate.json}; \file{Config/thresholds.json}; \file{Config/engineering_assumptions.json}.\par
}
\medskip
\textbf{Visualization limits.} Mission playback opens paused at year zero; Moon vectors come from the recording. Interpolation is display-only and monthly samples underresolve the Moon. Enlarged objects and schematic hardware are labelled. Live parameter edits require headless rerun/reload. The three-minute tour does not reach the coarse endpoint at its specified speed; use explicit endpoint seek. Viewer verification establishes display behaviour, not physics feasibility.

\end{document}

